\documentclass{article}
\usepackage[utf8]{inputenc}
\usepackage{hyperref}
\usepackage{listings}
\usepackage{xcolor}
\usepackage{enumitem}

\lstset{
    basicstyle=\ttfamily\small,
    breaklines=true,
    frame=single,
    backgroundcolor=\color{gray!5},
}

\title{Create detached PGP signatures for all files in a directory}
\author{Marcos de Carvalho}
\date{2026-04-28}

\begin{document}

\section*{Create detached PGP signatures for all files in a directory}
\textbf{Author:} Marcos de Carvalho \hfill \textbf{Date:} 2026-04-28

\noindent Iterates over every file in a given directory and creates a detached PGP signature (.sig) for each one using GnuPG.

\section*{Language}
bash

\section*{Category}
cryptography

\section*{Command}
\begin{lstlisting}
bash -c 'dir=$1; shift; for f in "$dir"/*; do [ -f "$f" ] && gpg --detach-sign "$f"; done' _
\end{lstlisting}

\section*{Explanation}
This oneliner uses a bash -c subshell to accept the target directory as a trailing argument. The variable dir receives \$1 (the directory path), then shift removes it so remaining arguments are available via \$@. A for loop iterates over all entries matching "\$dir"/*, and [ -f "\$f" ] ensures only regular files are processed (skipping subdirectories, symlinks, etc.). For each file, gpg --detach-sign creates a separate .sig file containing the PGP signature. The \_ placeholder serves as \$0 inside the subshell, ensuring \$1 maps to the first real argument.

\section*{Tags}
gpg, pgp, cryptography, signature, detach-sign, bulk, verification

\section*{Dependencies}
gnupg


\section*{Arguments}
\begin{enumerate}
  \item \textbf{DIR} (Optional): Path to the directory whose files will be signed. \\ Default: .
\end{enumerate}

\section*{Examples}
\begin{enumerate}
  \item \texttt{bash -c 'dir=\$1; shift; for f in "\$dir"/*; do [ -f "\$f" ] \&\& gpg --detach-sign "\$f"; done' \_ /path/to/releases} - Sign all files in /path/to/releases
  \item \texttt{bash -c 'dir=\$1; shift; for f in "\$dir"/*; do [ -f "\$f" ] \&\& gpg --detach-sign "\$f"; done' \_ .} - Sign all files in the current directory
  \item \texttt{alias bsign="bash -c 'dir=\textbackslash\{\}\$1; shift; for f in \textbackslash\{\}"\textbackslash\{\}\$dir\textbackslash\{\}"/*; do [ -f \textbackslash\{\}"\textbackslash\{\}\$f\textbackslash\{\}" ] \&\& gpg --detach-sign \textbackslash\{\}"\textbackslash\{\}\$f\textbackslash\{\}"; done' \_"} - Create an alias named 'bsign' for easy bulk signing
  \item \texttt{bsign /path/to/releases} - Use the alias to sign all files in a directory
\end{enumerate}

\section*{Output}
\begin{lstlisting}
file1.tar.gz.sig
file2.iso.sig
file3.deb.sig
\end{lstlisting}

\section*{Notes}
\begin{itemize}
  \item Each original file gets a corresponding .sig file in the same directory
  \item GPG will prompt for passphrase if the private key is passphrase-protected
  \item Uses the default signing key; set GPGKEY env var or use --default-key to specify another
  \item The [ -f "\$f" ] check skips subdirectories and non-regular files
  \item Hidden files (dotfiles) are not matched by the glob *; use .* as well if needed
  \item For large numbers of files, consider using gpg --batch --yes to skip confirmation prompts
\end{itemize}

\section*{Warnings}
\begin{itemize}
  \item Existing .sig files will be silently overwritten without confirmation
  \item Ensure you have sufficient disk space for the signature files
  \item If GPG cannot access your private key, the loop will fail for every file
  \item Signatures are created in the same directory as the source files
\end{itemize}

\section*{See Also}
\begin{itemize}
  \item find-large-files-recursive
\end{itemize}

\section*{Status}
reviewed

\section*{Safety}
caution

\section*{Shell}
bash

\section*{Platforms}
linux, gnu-linux, freebsd, openbsd, netbsd

\section*{Created At}
2026-04-28

\section*{Updated At}
2026-04-28

\end{document}
